% Options for packages loaded elsewhere
\PassOptionsToPackage{unicode}{hyperref}
\PassOptionsToPackage{hyphens}{url}
\documentclass[
]{article}
\usepackage{xcolor}
\usepackage[margin=1in]{geometry}
\usepackage{amsmath,amssymb}
\setcounter{secnumdepth}{-\maxdimen} % remove section numbering
\usepackage{iftex}
\ifPDFTeX
  \usepackage[T1]{fontenc}
  \usepackage[utf8]{inputenc}
  \usepackage{textcomp} % provide euro and other symbols
\else % if luatex or xetex
  \usepackage{unicode-math} % this also loads fontspec
  \defaultfontfeatures{Scale=MatchLowercase}
  \defaultfontfeatures[\rmfamily]{Ligatures=TeX,Scale=1}
\fi
\usepackage{lmodern}
\ifPDFTeX\else
  % xetex/luatex font selection
\fi
% Use upquote if available, for straight quotes in verbatim environments
\IfFileExists{upquote.sty}{\usepackage{upquote}}{}
\IfFileExists{microtype.sty}{% use microtype if available
  \usepackage[]{microtype}
  \UseMicrotypeSet[protrusion]{basicmath} % disable protrusion for tt fonts
}{}
\makeatletter
\@ifundefined{KOMAClassName}{% if non-KOMA class
  \IfFileExists{parskip.sty}{%
    \usepackage{parskip}
  }{% else
    \setlength{\parindent}{0pt}
    \setlength{\parskip}{6pt plus 2pt minus 1pt}}
}{% if KOMA class
  \KOMAoptions{parskip=half}}
\makeatother
\usepackage{longtable,booktabs,array}
\newcounter{none} % for unnumbered tables
\usepackage{calc} % for calculating minipage widths
% Correct order of tables after \paragraph or \subparagraph
\usepackage{etoolbox}
\makeatletter
\patchcmd\longtable{\par}{\if@noskipsec\mbox{}\fi\par}{}{}
\makeatother
% Allow footnotes in longtable head/foot
\IfFileExists{footnotehyper.sty}{\usepackage{footnotehyper}}{\usepackage{footnote}}
\makesavenoteenv{longtable}
\setlength{\emergencystretch}{3em} % prevent overfull lines
\providecommand{\tightlist}{%
  \setlength{\itemsep}{0pt}\setlength{\parskip}{0pt}}
\usepackage{bookmark}
\IfFileExists{xurl.sty}{\usepackage{xurl}}{} % add URL line breaks if available
\urlstyle{same}
\hypersetup{
  pdftitle={Computational Overhead of FIPS 203/204/205 Post-Quantum Cryptography in a BFT Consensus Pipeline},
  hidelinks,
  pdfcreator={LaTeX via pandoc}}

\title{Computational Overhead of FIPS 203/204/205 Post-Quantum
Cryptography in a BFT Consensus Pipeline}
\author{}
\date{}

\begin{document}
\maketitle

\begin{abstract}

This paper reports the computational overhead of NIST post-quantum
cryptographic standards (ML-DSA-65, ML-KEM-768, SLH-DSA-128s) integrated
into a HotStuff-based BFT consensus implementation. We instrument the
complete block production pipeline of a 4-validator configuration ---
key generation, transaction signing, consensus voting, signature
verification, and parallel transaction execution --- running in a single
process without network transport, and compare classical (Ed25519)
against post-quantum (ML-DSA-65) configurations. This isolates
cryptographic cost from network cost. In-process throughput with
ML-DSA-65 reaches 10,260 TPS with 9.75 ms block latency on a
laptop-class CPU using unoptimized reference implementations. The
primary bottleneck is transaction signing (83.2\% of total time);
verification overhead is small (1.3x). We validate implementation
correctness against 117 official NIST ACVP test vectors with byte-exact
output matching. Results are reproducible from the open-source
implementation.

\end{abstract}

\section{1. Introduction}\label{introduction}

NIST finalized three post-quantum cryptographic standards in August
2024: FIPS 203 (ML-KEM), FIPS 204 (ML-DSA), and FIPS 205 (SLH-DSA). The
final NIST IR 8547 (October 2026) deprecates quantum-vulnerable
signature and key-establishment algorithms after 2030 and disallows them
after 2035, regardless of key size {[}8{]}. Resource estimates for
quantum attacks keep decreasing: Gidney estimated that RSA-2048 can be
factored with fewer than one million noisy qubits in under a week
{[}1{]}, and Google Quantum AI estimated that secp256k1 ECDLP can be
solved with fewer than 500,000 physical qubits in 18 to 23 minutes on a
superconducting architecture, of which 9 to 12 minutes follow a
precomputation phase {[}2{]}; earlier analyses show how strongly such
estimates depend on hardware assumptions {[}6{]}. Because migrating
long-lived systems takes years, it has to finish before such a machine
exists {[}5{]}{[}10{]}.

Prior work on PQC performance focuses mostly on isolated cryptographic
operations {[}3{]}{[}4{]}. Two recent works measure post-quantum
signatures inside blockchain systems. Yu et al.~{[}11{]} replace BLS
with ML-DSA-65 in a BFT protocol deployed on 50 nodes across 5 cloud
regions and report that latency roughly doubles, to about one second.
Santos et al.~{[}12{]} run more than 11 million transactions on an
EVM-compatible chain with cryptographically agile transactions
supporting ECDSA, Falcon-512 and ML-DSA, and report no measurable
overhead from the agile transaction format. BFT protocols amplify
cryptographic costs because each round requires O(n) signatures and up
to O(n\^{}2) verifications across validators.

This paper complements those network-level results with a per-phase
breakdown of where cryptographic cost goes inside a single validator's
pipeline: consensus voting, transaction signing and execution, plus
hybrid signatures, key exchange and block size.

\section{2. System Architecture}\label{system-architecture}

The system under test implements a DAG-based HotStuff BFT consensus
protocol with Delegated Proof-of-Stake validator selection. The
cryptographic module is isolated in a separate crate with a state
machine modeled on FIPS 140-3 (4 states, 3 transitions, power-up
self-tests); the module has not been validated by a CMVP laboratory.

\subsection{2.1 Cryptographic
Primitives}\label{cryptographic-primitives}

{\def\LTcaptype{none} % do not increment counter
\begin{longtable}[]{@{}
  >{\raggedright\arraybackslash}p{(\linewidth - 8\tabcolsep) * \real{0.2075}}
  >{\raggedright\arraybackslash}p{(\linewidth - 8\tabcolsep) * \real{0.1887}}
  >{\raggedright\arraybackslash}p{(\linewidth - 8\tabcolsep) * \real{0.1698}}
  >{\raggedright\arraybackslash}p{(\linewidth - 8\tabcolsep) * \real{0.2264}}
  >{\raggedright\arraybackslash}p{(\linewidth - 8\tabcolsep) * \real{0.2075}}@{}}
\toprule\noalign{}
\begin{minipage}[b]{\linewidth}\raggedright
Primitive
\end{minipage} & \begin{minipage}[b]{\linewidth}\raggedright
Standard
\end{minipage} & \begin{minipage}[b]{\linewidth}\raggedright
Purpose
\end{minipage} & \begin{minipage}[b]{\linewidth}\raggedright
Public Key
\end{minipage} & \begin{minipage}[b]{\linewidth}\raggedright
Signature
\end{minipage} \\
\midrule\noalign{}
\endhead
\bottomrule\noalign{}
\endlastfoot
ML-DSA-65 & FIPS 204 & Block signatures, BFT votes, certificates & 1,952
B & 3,309 B \\
SLH-DSA-SHAKE-128s & FIPS 205 & Backup signatures (hash-based) & 32 B &
7,856 B \\
ML-KEM-768 & FIPS 203 & TLS key exchange (hybrid with X25519) & 1,184 B
& ct: 1,088 B \\
Ed25519 & RFC 8032 & Legacy/hybrid classical component & 32 B & 64 B \\
SHA3-256 & FIPS 202 & Block hashing, content hashing & --- & 32 B \\
\end{longtable}
}

\subsection{2.2 Consensus Protocol}\label{consensus-protocol}

HotStuff BFT with 3 phases (Prepare, PreCommit, Commit). Quorum
threshold: 2f+1 out of 3f+1 validators. Each phase requires 2f+1 valid
signatures. A single consensus round produces 9 signatures (3 phases x 3
validators) and 9 verifications.

\subsection{2.3 Hybrid Mode}\label{hybrid-mode}

Following ANSSI {[}7{]} and BSI {[}9{]} recommendations, the system
supports dual signatures (classical + post-quantum) on signed envelopes.
Both signatures must verify for the envelope to be accepted. The hybrid
TLS layer uses X25519 + ML-KEM-768 via rustls-post-quantum.

\section{3. Methodology}\label{methodology}

\subsection{3.1 Test Environment}\label{test-environment}

\begin{itemize}
\tightlist
\item
  Hardware: Apple M-series, 8 cores, 16 GB RAM
\item
  Toolchain: Rust nightly, release profile (optimizations enabled)
\item
  Crypto library: PQClean (C reference implementations via FFI)
\item
  Measurement: \texttt{std::time::Instant} (monotonic clock)
\item
  Iterations: 20 blocks x 100 transactions per benchmark run
\item
  Validators: 4 (minimum BFT with f=1), instantiated in one process;
  messages are passed in memory with no network transport, serialization
  over sockets, or propagation delay
\end{itemize}

\subsection{3.2 Pipeline
Instrumentation}\label{pipeline-instrumentation}

The benchmark measures three pipeline phases independently:

\begin{enumerate}
\def\labelenumi{\arabic{enumi}.}
\tightlist
\item
  \textbf{BFT Consensus}: Key generation, vote signing (9 signs per
  round), vote verification (9 verifies per round), state machine
  transitions.
\item
  \textbf{Transaction Signing}: Each transaction signed by the submitter
  using the configured algorithm.
\item
  \textbf{Execution}: Parallel transaction execution against in-memory
  world state with MVCC conflict detection.
\end{enumerate}

Wall-clock time recorded per phase per block. Total time includes all
phases plus overhead (serialization, memory allocation).

\subsection{3.3 ACVP Validation}\label{acvp-validation}

Implementation correctness verified against official NIST ACVP-Server
test vectors (github.com/usnistgov/ACVP-Server). Results:

{\def\LTcaptype{none} % do not increment counter
\begin{longtable}[]{@{}llll@{}}
\toprule\noalign{}
Algorithm & Test & Vectors & Result \\
\midrule\noalign{}
\endhead
\bottomrule\noalign{}
\endlastfoot
ML-DSA-65 & keyGen & 25 & 25/25 byte-exact \\
ML-DSA-65 & sigGen (deterministic) & 30 & 30/30 byte-exact \\
ML-DSA-65 & sigVer (pure mode) & 12 & 12/12 correct \\
ML-KEM-768 & keyGen & 25 & 25/25 byte-exact \\
ML-KEM-768 & encapsulation & 25 & 25/25 byte-exact \\
\end{longtable}
}

\section{4. Results}\label{results}

\subsection{4.1 Isolated Cryptographic
Operations}\label{isolated-cryptographic-operations}

100 iterations per algorithm (SLH-DSA: 10), release build.

{\def\LTcaptype{none} % do not increment counter
\begin{longtable}[]{@{}lllll@{}}
\toprule\noalign{}
Operation & Ed25519 & ML-DSA-65 & SLH-DSA-128s & ML-DSA/Ed25519 \\
\midrule\noalign{}
\endhead
\bottomrule\noalign{}
\endlastfoot
KeyGen (us) & 9.2 & 39.1 & 456,048 & 4.2x \\
Sign (us) & 9.1 & 90.5 & 352,797 & 9.9x \\
Verify (us) & 19.6 & 25.5 & 352.5 & 1.3x \\
Public key (B) & 32 & 1,952 & 32 & 61.0x \\
Signature (B) & 64 & 3,309 & 7,856 & 51.7x \\
\end{longtable}
}

SLH-DSA-128s signing is 3,880x slower than ML-DSA-65, confirming its
role as emergency backup only.

\subsection{4.2 End-to-End Pipeline}\label{end-to-end-pipeline}

20 blocks, 100 transactions per block, 4 validators, real cryptography.

{\def\LTcaptype{none} % do not increment counter
\begin{longtable}[]{@{}llll@{}}
\toprule\noalign{}
Metric & Ed25519 & ML-DSA-65 & Overhead \\
\midrule\noalign{}
\endhead
\bottomrule\noalign{}
\endlastfoot
Throughput (TPS) & 64,561 & 10,260 & 6.3x \\
Block rate (blocks/s) & 646 & 103 & 6.3x \\
Block latency (ms) & 1.55 & 9.75 & 6.3x \\
\end{longtable}
}

\subsection{4.3 Time Breakdown}\label{time-breakdown}

{\def\LTcaptype{none} % do not increment counter
\begin{longtable}[]{@{}lllll@{}}
\toprule\noalign{}
Phase & Ed25519 (ms) & Ed25519 (\%) & ML-DSA-65 (ms) & ML-DSA-65 (\%) \\
\midrule\noalign{}
\endhead
\bottomrule\noalign{}
\endlastfoot
BFT consensus & 6.20 & 20.0 & 28.67 & 14.7 \\
TX signing & 20.73 & 66.9 & 162.22 & 83.2 \\
Execution & 3.62 & 11.7 & 3.53 & 1.8 \\
\end{longtable}
}

Transaction signing dominates in both configurations. The shift from
66.9\% to 83.2\% reflects the 9.9x signing overhead of ML-DSA-65 while
execution cost remains constant.

\subsection{4.4 Scaling Behavior}\label{scaling-behavior}

ML-DSA-65, 10 blocks per configuration.

{\def\LTcaptype{none} % do not increment counter
\begin{longtable}[]{@{}llll@{}}
\toprule\noalign{}
Txs/Block & TPS & Block Latency (ms) & Bottleneck \\
\midrule\noalign{}
\endhead
\bottomrule\noalign{}
\endlastfoot
10 & 4,442 & 2.25 & BFT \\
50 & 8,956 & 5.58 & Signing \\
100 & 10,041 & 9.96 & Signing \\
200 & 10,790 & 18.54 & Signing \\
500 & 11,014 & 45.40 & Signing \\
\end{longtable}
}

TPS plateaus around 11,000 as signing becomes the dominant cost. At low
transaction counts (\textless50), BFT consensus overhead dominates. The
crossover occurs near 30 transactions per block.

\subsection{4.5 Hybrid Signature
Overhead}\label{hybrid-signature-overhead}

Hybrid dual signature as recommended by ANSSI (Ed25519 primary +
ML-DSA-65 secondary), 50 iterations.

{\def\LTcaptype{none} % do not increment counter
\begin{longtable}[]{@{}llll@{}}
\toprule\noalign{}
Operation & Classical (us) & Hybrid (us) & Overhead \\
\midrule\noalign{}
\endhead
\bottomrule\noalign{}
\endlastfoot
Sign & 12.6 & 85.4 & 6.8x \\
Verify & 18.9 & 64.2 & 3.4x \\
Bandwidth per envelope & 96 B & 5,325 B & 55.5x \\
\end{longtable}
}

\subsection{4.6 Key Exchange (TLS)}\label{key-exchange-tls}

ML-KEM-768, 100 iterations.

{\def\LTcaptype{none} % do not increment counter
\begin{longtable}[]{@{}ll@{}}
\toprule\noalign{}
Operation & Latency (us) \\
\midrule\noalign{}
\endhead
\bottomrule\noalign{}
\endlastfoot
KeyGen & 9.8 \\
Encapsulate & 9.7 \\
Decapsulate & 8.9 \\
Handshake overhead & +2,208 B \\
\end{longtable}
}

\subsection{4.7 Block Size Impact}\label{block-size-impact}

Estimated block size with N endorsements (signature + payload hash +
public key per endorsement).

{\def\LTcaptype{none} % do not increment counter
\begin{longtable}[]{@{}llll@{}}
\toprule\noalign{}
Endorsements & Ed25519 (KB) & ML-DSA-65 (KB) & Ratio \\
\midrule\noalign{}
\endhead
\bottomrule\noalign{}
\endlastfoot
1 & 0.4 & 8.5 & 23.8x \\
5 & 1.0 & 29.2 & 29.7x \\
10 & 1.8 & 55.1 & 31.2x \\
20 & 3.3 & 106.8 & 32.1x \\
\end{longtable}
}

\section{5. Discussion}\label{discussion}

\subsection{5.1 Verification is Not the
Bottleneck}\label{verification-is-not-the-bottleneck}

The 1.3x verification overhead of ML-DSA-65 contradicts the assumption
that PQC significantly impacts consensus validation. In HotStuff BFT,
validators verify O(n) signatures per phase --- at 25.5 us per verify, a
100-validator network adds only 2.55 ms per phase, well within typical
network round-trip times.

\subsection{5.2 Signing Dominates}\label{signing-dominates}

Transaction signing consumes 83.2\% of pipeline time with ML-DSA-65.
This cost is per-transaction and scales linearly. Parallelizing
transaction signing across available cores is the most direct
optimization path --- the current implementation signs sequentially.

\subsection{5.3 Bandwidth
Considerations}\label{bandwidth-considerations}

ML-DSA-65 signatures (3,309 B) are 51.7x larger than Ed25519 (64 B). For
a block with 100 transactions and 5 endorsements each, the bandwidth
difference is approximately 1.5 MB (ML-DSA) vs 30 KB (Ed25519). At 103
blocks/second, this produces approximately 155 MB/s of signature data.
This is within the capacity of modern datacenter networks but may
constrain geographic distribution over slower links.

\subsection{5.4 Comparison with Existing Work and
Systems}\label{comparison-with-existing-work-and-systems}

Yu et al.~{[}11{]} measure ML-DSA-65 in a geo-distributed BFT deployment
and find roughly 2x latency, an effect dominated by network and message
size. Our in-process measurements show the purely computational side:
verification costs 1.3x and signing 9.9x Ed25519, so in a networked
deployment the bandwidth of 3,309-byte signatures (Section 5.3), not
verification time, is the more likely source of added latency. Santos et
al.~{[}12{]} show that supporting several signature schemes at the
transaction-format level adds no measurable overhead, which is
consistent with our hybrid-envelope design.

Among deployed systems, Algorand executed its first Falcon-signed
transaction on mainnet in November 2025 and enabled native Falcon-1024
accounts in August 2026 {[}13{]}. QRL 2.0 uses ML-DSA-87 for accounts
and validators and is on public testnet as of 2026 {[}14{]}. As of
November 2026, neither Bitcoin nor Ethereum has deployed a post-quantum
signature scheme on mainnet; Bitcoin's BIP-360 removes the Taproot
key-path exposure without introducing a post-quantum scheme {[}15{]},
and Ethereum's post-quantum migration targets 2029 {[}16{]}.

\subsection{5.5 Limitations}\label{limitations}

All validators run in one process, so the reported throughput excludes
network latency, bandwidth limits and real message serialization; it is
an upper bound on what the cryptographic pipeline allows, not a network
throughput figure. Measurements use PQClean reference implementations
without AVX2 or NEON optimizations, on a single laptop-class machine;
optimized implementations would reduce ML-DSA signing cost. Transaction
signing is sequential, which overstates its share of pipeline time on
multi-core hardware. A networked multi-machine evaluation is left for
future work.

\section{6. Reproducibility}\label{reproducibility}

All benchmarks are executed via:

\begin{verbatim}
cargo test --release --test e2e_throughput -- --nocapture
cargo test --release --test pqc_benchmark -- --nocapture
\end{verbatim}

Source code, test vectors, and benchmark harness are available in the
repository. ACVP validation vectors sourced from
github.com/usnistgov/ACVP-Server.

\section{References}\label{references}

{[}1{]} C. Gidney, ``How to factor 2048 bit RSA integers with less than
a million noisy qubits,'' arXiv:2505.15917, 2025.

{[}2{]} Google Quantum AI, ``Securing Elliptic Curve Cryptocurrencies
against Quantum Vulnerabilities: Resource Estimates and Mitigations,''
arXiv:2603.28846, 2026.

{[}3{]} M. J. Kannwischer et al., ``Improving software quality in
cryptography standardization projects,'' in IEEE EuroS\&P Workshops,
2022.

{[}4{]} NIST, ``Post-Quantum Cryptography: FIPS 203, 204, 205,''
National Institute of Standards and Technology, 2024.

{[}5{]} D. J. Bernstein and T. Lange, ``Post-quantum cryptography,''
Nature, vol.~549, pp.~188-194, 2017.

{[}6{]} M. Webber et al., ``The impact of hardware specifications on
reaching quantum advantage in the fault tolerant regime,'' AVS Quantum
Science, 2022.

{[}7{]} ANSSI, ``Avis relatif a la migration vers la cryptographie
post-quantique,'' Agence nationale de la securite des systemes
d'information, 2024.

{[}8{]} NIST, ``Transition to Post-Quantum Cryptography Standards,''
NIST IR 8547 (final), October 2026.

{[}9{]} BSI, ``Kryptographische Verfahren: Empfehlungen und
Schlussellangen,'' BSI TR-02102-1, 2024.

{[}10{]} M. Mosca, ``Cybersecurity in an era with quantum computers:
will we be ready?'' IEEE Security \& Privacy, vol.~16, no. 5, pp.~38-41,
2018.

{[}11{]} Q. Yu, J. Villacis, G. Losa, Z. Xiang and X. Wang, ``Simple-IT:
Practical Low-Latency Signature-Free BFT Consensus,'' arXiv:2606.14404,
2026.

{[}12{]} M. B. Santos, D. Ferrin, R. Kahat and M. Lodder, ``Post-Quantum
Blockchains with Agility in Mind,'' IACR Cryptology ePrint Archive,
2026/609, 2026.

{[}13{]} Algorand Foundation, ``Technical Brief: Quantum-resistant
transactions on Algorand with Falcon signatures,'' 2025; ``2025 on
Algorand: Roadmap progress,'' 2025.

{[}14{]} The QRL Foundation, ``QRL 2.0 (Project Zond),''
https://www.theqrl.org/2/, accessed November 2026.

{[}15{]} BIP-360, ``Pay to Merkle Root (P2MR),'' Bitcoin Improvement
Proposals, 2026.

{[}16{]} Ethereum Foundation, post-quantum roadmap, March 2026.

\end{document}
